\chapter{Eksperyment i wyniki}\label{chap:eksperyment-i-wyniki}

Rozdział dokumentuje przebieg eksperymentu porównawczego oraz surowe wyniki
liczbowe. Ramy metodyczne i~kryteria oceny opisano w~\cref{chap:metodologia},
a~szczegóły implementacji -- w~\cref{chap:implementacja}. Niniejszy rozdział
ogranicza się do prezentacji danych i~wyników -- interpretację i~wnioski
omówiono w~rozdziale następnym.

\section{Przebieg wykonania}\label{sec:przebieg-wykonania}

Eksperyment zrealizowano w~ramach jednego uruchomienia obejmującego pełną
macierz \(90 \times 6 = 540\) par pytanie--architektura. Poniższe podsekcje
dokumentują czas trwania, kompletność wyników oraz zestaw wygenerowanych
artefaktów.

\subsection{Czas trwania}\label{subsec:czas-trwania}

Uruchomienie przeprowadzono 13~marca 2026~r. pod identyfikatorem
\texttt{20260313\_200511}. Komponent wykonawczy przekazał 540~par do puli
32~wątków roboczych; ostatni rekord zapisano o~\texttt{21:52~UTC}, co odpowiada
rzeczywistemu czasowi wykonania \mbox{1~h~47~min}. Łączny czas sekwencyjny --
rozumiany jako suma opóźnień wszystkich uruchomień -- wynosi \mbox{62\,698~s}
(\(\approx\)17{,}4~h). Faza ewaluacji jakościowej w~systemie \texttt{Phoenix}
zakończyła się o~\texttt{22:55~UTC}.

Tabela~\ref{tab:czas-wykonania} przedstawia rozkład czasu sekwencyjnego według
architektur wraz z~udziałem procentowym oraz średnim czasem pojedynczego
uruchomienia (z~odchyleniem standardowym). Architektura ReAct wykazała
największy udział w~łącznym czasie sekwencyjnym (26,1\%), natomiast
planer--wykonawca -- najmniejszy (12,1\%). Stanowi to interesujący paradoks
wydajnościowy: mimo dodatkowego narzutu czasowego na generowanie planu,
agent ten okazał się najszybszy w~ogólnej egzekucji (średnio 84{,}3~s),
wyprzedzając pod tym względem nawet najprostszą architekturę deterministyczną
(93{,}0~s). Wynika to z~faktu, że trafnie ułożony plan pozwolił na redukcję
liczby zbędnych wywołań narzędzi i~kosztownych operacji pobierania danych.

\begin{table}[H]
  \caption{Czas wykonania według architektury}\label{tab:czas-wykonania}
  \centering
  \small
  \renewcommand{\arraystretch}{1.3}
  \begin{tabular}{@{} l r r r r @{}}
    \toprule
    \textbf{Architektura}
    & \textbf{Uruchomienia}
    & \textbf{Suma opóźnień [s]}
    & \textbf{Udział [\%]}
    & \textbf{Średnia \(\pm\) odchylenie [s]}                                                                          \\
    \midrule
    Deterministyczna    & 90                                      & 8\,372           & 13{,}4           & 93{,}0 (\(\pm\) 17{,}5)          \\
    Planer--wykonawca   & 90                                      & 7\,583           & 12{,}1           & \textbf{84{,}3 (\(\pm\) 17{,}1)} \\
    Router--specjalista & 90                                      & 8\,902           & 14{,}2           & 98{,}9 (\(\pm\) 18{,}3)          \\
    Tablicowa           & 90                                      & 8\,229           & 13{,}1           & 91{,}4 (\(\pm\) 19{,}5)          \\
    Hierarchiczna       & 90                                      & 13\,239          & 21{,}1           & 147{,}1 (\(\pm\) 68{,}2)         \\
    ReAct               & 90                                      & \textbf{16\,374} & \textbf{26{,}1}  & 181{,}9 (\(\pm\) 44{,}9)         \\
    \midrule
    \textbf{Łącznie}    & \textbf{540}                            & \textbf{62\,698} & \textbf{100{,}0} & --                               \\
    \bottomrule
  \end{tabular}
  \zrodlo{opracowanie własne}
\end{table}

\subsection{Kompletność i~błędy wykonania}\label{subsec:kompletnosc-i-bledy}

Spośród 540~uruchomień 539 zakończyło się pomyślnie. Jedyny błąd odnotowano
w~architekturze hierarchicznej dla pytania dotyczącego zakresu autocasco na
terytorium Białorusi. Bezpośrednią przyczyną był wyjątek
\texttt{openai.BadRequestError~400} zwrócony przez interfejs \texttt{OpenAI
API}: węzeł dekompozytora wygenerował podpytanie w~mieszanym formacie
polsko-angielskim, co uniemożliwiło poprawną serializację odpowiedzi. Wyjątek
przechwycono i~zapisano jako rekord z~flagą błędu; zgodnie z~przyjętą procedurą
nie ponowiono wykonania zadania. Wskaźnik błędów wynosi \(1/540 \approx
0{,}19\%\) dla całego eksperymentu i~\(1{,}11\%\) dla architektury
hierarchicznej.

Żadne uruchomienie nie przekroczyło limitu czasowego 120~s. Nie odnotowano również
osiągnięcia limitu iteracji w~architekturach iteracyjnych; maksymalna liczba iteracji
architektury ReAct wyniosła~11 (przy progu równym 15).

\subsection{Artefakty wynikowe}\label{subsec:artefakty-wynikowe}

Wyniki uruchomienia zapisano w~katalogu \texttt{code/data/evaluation/results/}.
Tabela~\ref{tab:artefakty-wynikowe} zestawia pliki wyjściowe z~ich rozmiarami
i~zawartością; gwiazdka zastępuje pełny identyfikator uruchomienia
\texttt{20260313\_200511}. Największy rozmiar ma plik \texttt{.csv} z~surowymi
wynikami, natomiast plik zbiorczego podsumowania zajmuje poniżej 1~KB.

\begin{table}[H]
  \caption{Pliki wynikowe eksperymentu}\label{tab:artefakty-wynikowe}
  \centering
  \small
  \renewcommand{\arraystretch}{1.3}
  \begin{tabularx}{\textwidth}{@{} l r L @{}}
    \toprule
    \textbf{Plik}
    & \textbf{Rozmiar}
    & \textbf{Zawartość}                                     \\
    \midrule
    \texttt{runner\_results\_*.jsonl}
    & 16\,006~KB
    & 540~rekordów surowych w~formacie strumieniowym         \\
    \texttt{runner\_results\_*.csv}
    & 17\,131~KB
    & spłaszczona tabela 540~wierszy i~32~kolumn             \\
    \texttt{runner\_summary\_*.csv}
    & 2~KB
    & wartości zagregowane według architektury               \\
    \texttt{overall\_answer\_metrics\_*.csv}
    & 1\,988~KB
    & 539~ocen jakości odpowiedzi z~systemu \texttt{Phoenix} \\
    \texttt{overall\_document\_relevance\_*.csv}
    & 2\,216~KB
    & 1\,786~ocen trafności pojedynczych cytowań             \\
    \texttt{overall\_summary\_*.csv}
    & \({<}\)1~KB
    & zbiorcze miary jakości według architektury             \\
    \bottomrule
  \end{tabularx}
  \zrodlo{opracowanie własne}
\end{table}

\section{Wyniki operacyjne}\label{sec:wyniki-operacyjne}

W sekcji zestawiono średnie wartości miar opisujących zasobochłonność
uruchomień: zużycie tokenów i~koszt, liczbę wywołań narzędzi i~iteracji oraz
liczbę pobranych fragmentów i~cytowań. Wszystkie wartości w~tabelach podano
w~formacie średnia \(\pm\) odchylenie standardowe; podstawę obliczeń stanowi
90~uruchomień danej architektury, z~wyjątkiem hierarchicznej (89~uruchomień).

\subsection{Zużycie tokenów i~koszt}\label{subsec:zuzycie-tokenow-i-koszt}

Tabela~\ref{tab:tokeny-i-koszt} podaje średnie zużycie tokenów wejściowych
i~wyjściowych oraz średni koszt pojedynczego uruchomienia, oszacowany na
podstawie cennika \texttt{OpenAI API} obowiązującego w~marcu 2026~r.
Architektury jednoagentowe wykazują niewielki rozrzut: średnie zużycie tokenów
wejściowych mieści się w~przedziale 2\,988--3\,219, a~odchylenie standardowe
nie przekracza 260~tokenów. Architektury iteracyjne wykazują znaczną zmienność
-- średnia hierarchicznej wynosi 450\,017~tokenów wejściowych przy odchyleniu
standardowym 1\,080\,756, a~architektury ReAct 391\,935 przy odchyleniu
419\,373. Wartości odchylenia standardowego tak drastycznie przekraczające średnią
wynikają z~silnie prawostronnie skośnego rozkładu zużycia tokenów: w~większości
uruchomień agenci zużywają stosunkowo niewielką liczbę tokenów, jednak pojedyncze
przebiegi, w~których agent ulega zapętleniu lub wielokrotnie ponawia błędne
wywołania, konsumują ekstremalnie duże zasoby, co diametralnie zawyża średnią
i~wariancję. Średni koszt uruchomienia wynosi od 0{,}00037~USD (architektura
tablicowa) do 0{,}04320~USD (ReAct).

\begin{table}[H]
  \caption{Tokeny i~koszt na uruchomienie}\label{tab:tokeny-i-koszt}
  \centering
  \small
  \renewcommand{\arraystretch}{1.3}
  \begin{tabular}{@{} l r r r @{}}
    \toprule
    \textbf{Architektura}
    & \textbf{Tokeny wejściowe}      & \textbf{Tokeny wyjściowe}   & \textbf{Koszt [USD]}                  \\
    \midrule
    Deterministyczna    & \textbf{2\,988 (\(\pm\) 230)}  & 528 (\(\pm\) 160)           & 0{,}00038 (\(\pm\) 0{,}00007)          \\
    Planer--wykonawca   & 3\,219 (\(\pm\) 260)           & 580 (\(\pm\) 159)           & 0{,}00040 (\(\pm\) 0{,}00007)          \\
    Router--specjalista & 3\,199 (\(\pm\) 246)           & 519 (\(\pm\) 151)           & 0{,}00039 (\(\pm\) 0{,}00007)          \\
    Tablicowa           & 3\,004 (\(\pm\) 215)           & \textbf{498 (\(\pm\) 144)}  & \textbf{0{,}00037 (\(\pm\) 0{,}00006)} \\
    Hierarchiczna       & 450\,017 (\(\pm\) 1\,080\,756) & 115\,761 (\(\pm\) 303\,962) & 0{,}07109 (\(\pm\) 0{,}18099)          \\
    ReAct               & 391\,935 (\(\pm\) 419\,373)    & 58\,548 (\(\pm\) 66\,709)   & 0{,}04320 (\(\pm\) 0{,}04758)          \\
    \bottomrule
  \end{tabular}
  \zrodlo{opracowanie własne}
\end{table}

\subsection{Wywołania narzędzi i~iteracje}\label{subsec:wywolania-i-iteracje}

Tabela~\ref{tab:wywolania-i-iteracje} podaje średnią liczbę wywołań narzędzi
i~iteracji na uruchomienie wraz z~odchyleniem standardowym. Pojęcie iteracji
zależy od topologii: dla planera--wykonawcy oznacza długość wygenerowanego
planu, dla tablicowej -- liczbę cykli dyspozytora, dla hierarchicznej -- liczbę
przetworzonych podpytań, dla ReAct -- liczbę zarejestrowanych obserwacji.
Architektury deterministyczna i~router--specjalista nie posługują się pojęciem
iteracji. Architektura hierarchiczna wykonuje średnio 1\,066 wywołań narzędzi
na uruchomienie przy odchyleniu 2\,568, a~architektura ReAct 295{,}7 przy
odchyleniu 294{,}1. Planer--wykonawca wykazuje najniższą liczbę wywołań
narzędzi przy najmniejszym odchyleniu (średnio 7{,}2 przy odchyleniu 0{,}8).

\begin{table}[H]
  \caption{Liczba wywołań narzędzi i~iteracji na uruchomienie}\label{tab:wywolania-i-iteracje}
  \centering
  \small
  \renewcommand{\arraystretch}{1.3}
  \begin{tabular}{@{} l r r @{}}
    \toprule
    \textbf{Architektura}
    & \textbf{Wywołania narzędzi}     & \textbf{Iteracje}              \\
    \midrule
    Deterministyczna    & 10{,}0 (\(\pm\) 0{,}0)          & --                             \\
    Planer--wykonawca   & \textbf{7{,}2 (\(\pm\) 0{,}8)}  & \textbf{4{,}2 (\(\pm\) 0{,}8)} \\
    Router--specjalista & 9{,}8 (\(\pm\) 0{,}6)           & --                             \\
    Tablicowa           & 10{,}0 (\(\pm\) 0{,}0)          & 8{,}0 (\(\pm\) 0{,}0)          \\
    Hierarchiczna       & 1\,066{,}0 (\(\pm\) 2\,568{,}0) & 10{,}6 (\(\pm\) 4{,}8)         \\
    ReAct               & 295{,}7 (\(\pm\) 294{,}1)       & 7{,}7 (\(\pm\) 1{,}4)          \\
    \bottomrule
  \end{tabular}
  \zrodlo{opracowanie własne}
\end{table}

\subsection{Fragmenty z~wyszukiwania i~cytowania}\label{subsec:fragmenty-i-cytowania}

Tabela~\ref{tab:fragmenty-i-cytowania} przedstawia średnie liczby fragmentów
pobranych z~bazy \texttt{Qdrant} oraz cytowań w~odpowiedzi końcowej wraz
z~odchyleniami standardowymi. Kolumna \textit{Suma} podaje łączną liczbę
fragmentów pobranych w~obrębie jednego uruchomienia -- w~architekturach
iteracyjnych obejmuje wszystkie wywołania pętli roboczej lub obserwacyjnej.
Średnia łączna liczba fragmentów dla architektury hierarchicznej wynosi 1\,772
przy odchyleniu 4\,300, a~dla ReAct 154{,}5 przy odchyleniu 131{,}1.
Architektury jednoagentowe pobierają 5 lub 15~fragmentów na uruchomienie przy
zerowym lub minimalnym odchyleniu standardowym. W~architekturze ReAct liczba
cytowań w~odpowiedzi końcowej jest niższa niż w~pozostałych wariantach (średnio
  w~przedziale 3{,}29--3{,}82.

  Szczególną anomalię odnotowano w~przypadku architektury planera--wykonawcy.
  Zgodnie z~tabelą~\ref{tab:fragmenty-i-cytowania}, po etapie oceny trafności
  w~stanie agenta pozostaje średnio zaledwie 0{,}31 fragmentu, co świadczy
  o~niezwykle rygorystycznym odrzucaniu pobranego kontekstu przez sędziego.
  Paradoksalnie, odpowiedź końcowa generowana przez tę architekturę posiada
  średnio aż 3{,}38 cytowania. Wskazuje to na iluzję filtrowania dowodów:
  moduł syntezy w~przypadku planera zignorował wyniki oceny trafności
  i~samodzielnie uwzględnił fragmenty (bądź wyhalucynował ich znaczniki),
  które na etapie oceny zostały z~kontekstu usunięte.

  \begin{table}[H]
    \caption{Liczba fragmentów i~cytowań}\label{tab:fragmenty-i-cytowania}
    \centering
    \small
    \renewcommand{\arraystretch}{1.3}
    \begin{tabular}{@{} l r r r r @{}}
      \toprule
      \textbf{Architektura}
      & \textbf{Na zapytanie}            & \textbf{Suma}                   & \textbf{Po ocenie trafności}    & \textbf{Cytowania}               \\
      \midrule
      Deterministyczna    & \textbf{4{,}28 (\(\pm\) 1{,}00)} & 15{,}0 (\(\pm\) 0{,}0)          & 4{,}28 (\(\pm\) 1{,}00)         & 3{,}72 (\(\pm\) 0{,}67)          \\
      Planer--wykonawca   & 4{,}98 (\(\pm\) 0{,}15)          & \textbf{5{,}0 (\(\pm\) 0{,}0)}  & \textbf{0{,}31 (\(\pm\) 1{,}17)} & 3{,}38 (\(\pm\) 0{,}77)          \\
      Router--specjalista & 4{,}33 (\(\pm\) 0{,}99)          & 14{,}4 (\(\pm\) 2{,}0)          & 4{,}33 (\(\pm\) 0{,}99)         & 3{,}67 (\(\pm\) 0{,}65)          \\
      Tablicowa           & 4{,}30 (\(\pm\) 0{,}90)          & 15{,}0 (\(\pm\) 0{,}0)          & 4{,}30 (\(\pm\) 0{,}90)         & \textbf{3{,}74 (\(\pm\) 0{,}53)} \\
      Hierarchiczna       & 4{,}70 (\(\pm\) 1{,}48)          & 1\,771{,}7 (\(\pm\) 4\,299{,}8) & 4{,}70 (\(\pm\) 1{,}48)         & 3{,}74 (\(\pm\) 0{,}71)          \\
      ReAct               & 4{,}39 (\(\pm\) 1{,}07)          & 154{,}5 (\(\pm\) 131{,}1)       & 4{,}22 (\(\pm\) 1{,}32)         & 1{,}59 (\(\pm\) 1{,}69)          \\
      \bottomrule
    \end{tabular}
    \zrodlo{opracowanie własne}
  \end{table}

  \section{Wyniki jakościowe}\label{sec:wyniki-jakosciowe}

  Poniższe podsekcje przedstawiają wartości czterech miar jakości w~rozbiciu na
  architekturę, poziom trudności pytania, kategorię produktową i~ubezpieczyciela.
  Liczebności podzbiorów wynoszą: 9~pytań bardzo prostych, 27~prostych,
  27~złożonych i~27~bardzo złożonych; 30~pytań na każdą kategorię produktową
  i~30~na każdego ubezpieczyciela. W~architekturze hierarchicznej poziom złożony
  obejmuje 26~pytań ze względu na pojedynczy błąd wykonania.

  \subsection{Wierność}\label{subsec:wyniki-wiernosc}

  Tabele~\ref{tab:wiernosc-trudnosc}--\ref{tab:wiernosc-ubezpieczyciel}
  przedstawiają średnią wierność odpowiedzi wraz z~odchyleniem standardowym
  w~rozbiciu na poziom trudności, kategorię produktową i~ubezpieczyciela.
  Najwyższą wartość łączną osiąga router--specjalista (0{,}811), najniższą ReAct
  (0{,}611). Dla wszystkich architektur wierność spada wraz ze wzrostem poziomu
  trudności pytania. W~rozbiciu według kategorii produktowej i~ubezpieczyciela
  rozpiętość wyników między architekturami pozostaje zbliżona do rozpiętości
  w~ujęciu łącznym, z~jednym silnym wyjątkiem. Architektura ReAct odnotowała
  katastrofalny spadek wierności w~kategorii ubezpieczeń mieszkaniowych,
  uzyskując zaledwie 0{,}433 (wobec 0{,}667 dla ubezpieczeń komunikacyjnych
  i~0{,}733 dla turystycznych). Wynik bliski losowości wskazuje na wybitną
  podatność pętli obserwacyjnej ReAct na dryf semantyczny w~kategoriach
  o~szerszym i~bardziej podatnym na wieloznaczność zakresie pojęciowym.

  \begin{table}[H]
    \caption{Wierność według poziomu trudności}\label{tab:wiernosc-trudnosc}
    \centering
    \footnotesize
    \renewcommand{\arraystretch}{1.3}
    \begin{tabular}{@{} l c c c c c @{}}
      \toprule
      \textbf{Architektura}
      & \textbf{B.~proste}
      & \textbf{Proste}
      & \textbf{Złożone}
      & \textbf{B.~złożone}
      & \textbf{Łącznie}                                                                                                                                                                  \\
      \midrule
      Deterministyczna    & \textbf{0{,}889 (\(\pm\) 0{,}105)} & \textbf{0{,}852 (\(\pm\) 0{,}068)} & 0{,}741 (\(\pm\) 0{,}084)          & 0{,}704 (\(\pm\) 0{,}088)          & 0{,}778 (\(\pm\) 0{,}044)          \\
      Planer--wykonawca   & 0{,}889 (\(\pm\) 0{,}105)          & 0{,}741 (\(\pm\) 0{,}084)          & \textbf{0{,}852 (\(\pm\) 0{,}068)} & 0{,}741 (\(\pm\) 0{,}084)          & 0{,}789 (\(\pm\) 0{,}043)          \\
      Router--specjalista & 0{,}889 (\(\pm\) 0{,}105)          & 0{,}852 (\(\pm\) 0{,}068)          & 0{,}741 (\(\pm\) 0{,}084)          & \textbf{0{,}815 (\(\pm\) 0{,}075)} & \textbf{0{,}811 (\(\pm\) 0{,}041)} \\
      Tablicowa           & 0{,}778 (\(\pm\) 0{,}139)          & 0{,}852 (\(\pm\) 0{,}068)          & 0{,}704 (\(\pm\) 0{,}088)          & 0{,}815 (\(\pm\) 0{,}075)          & 0{,}789 (\(\pm\) 0{,}043)          \\
      Hierarchiczna       & 0{,}556 (\(\pm\) 0{,}166)          & 0{,}815 (\(\pm\) 0{,}075)          & 0{,}731 (\(\pm\) 0{,}087)          & 0{,}593 (\(\pm\) 0{,}095)          & 0{,}697 (\(\pm\) 0{,}049)          \\
      ReAct               & 0{,}556 (\(\pm\) 0{,}166)          & 0{,}667 (\(\pm\) 0{,}091)          & 0{,}519 (\(\pm\) 0{,}096)          & 0{,}667 (\(\pm\) 0{,}091)          & 0{,}611 (\(\pm\) 0{,}051)          \\
      \bottomrule
    \end{tabular}
    \zrodlo{opracowanie własne}
  \end{table}

  \begin{table}[H]
    \caption{Wierność według kategorii produktowej}\label{tab:wiernosc-kategoria}
    \centering
    \small
    \renewcommand{\arraystretch}{1.3}
    \begin{tabular}{@{} l c c c c @{}}
      \toprule
      \textbf{Architektura}
      & \textbf{Auto}
      & \textbf{Mieszkanie}
      & \textbf{Podróż}
      & \textbf{Łącznie}                                                                                                                              \\
      \midrule
      Deterministyczna    & 0{,}800 (\(\pm\) 0{,}073)          & 0{,}800 (\(\pm\) 0{,}073)          & 0{,}733 (\(\pm\) 0{,}081)          & 0{,}778 (\(\pm\) 0{,}044)          \\
      Planer--wykonawca   & 0{,}800 (\(\pm\) 0{,}073)          & \textbf{0{,}833 (\(\pm\) 0{,}068)} & 0{,}733 (\(\pm\) 0{,}081)          & 0{,}789 (\(\pm\) 0{,}043)          \\
      Router--specjalista & 0{,}800 (\(\pm\) 0{,}073)          & 0{,}833 (\(\pm\) 0{,}068)          & \textbf{0{,}800 (\(\pm\) 0{,}073)} & \textbf{0{,}811 (\(\pm\) 0{,}041)} \\
      Tablicowa           & \textbf{0{,}833 (\(\pm\) 0{,}068)} & 0{,}800 (\(\pm\) 0{,}073)          & 0{,}733 (\(\pm\) 0{,}081)          & 0{,}789 (\(\pm\) 0{,}043)          \\
      Hierarchiczna       & 0{,}586 (\(\pm\) 0{,}091)          & 0{,}700 (\(\pm\) 0{,}084)          & 0{,}800 (\(\pm\) 0{,}073)          & 0{,}697 (\(\pm\) 0{,}049)          \\
      ReAct               & 0{,}667 (\(\pm\) 0{,}086)          & 0{,}433 (\(\pm\) 0{,}090)          & 0{,}733 (\(\pm\) 0{,}081)          & 0{,}611 (\(\pm\) 0{,}051)          \\
      \bottomrule
    \end{tabular}
    \zrodlo{opracowanie własne}
  \end{table}

  \begin{table}[H]
    \caption{Wierność według ubezpieczyciela}\label{tab:wiernosc-ubezpieczyciel}
    \centering
    \small
    \renewcommand{\arraystretch}{1.3}
    \begin{tabular}{@{} l c c c c @{}}
      \toprule
      \textbf{Architektura}
      & \textbf{Hestia}
      & \textbf{PZU}
      & \textbf{Warta}
      & \textbf{Łącznie}                                                                                                                              \\
      \midrule
      Deterministyczna    & 0{,}700 (\(\pm\) 0{,}084)          & 0{,}800 (\(\pm\) 0{,}073)          & 0{,}833 (\(\pm\) 0{,}068)          & 0{,}778 (\(\pm\) 0{,}044)          \\
      Planer--wykonawca   & 0{,}733 (\(\pm\) 0{,}081)          & 0{,}800 (\(\pm\) 0{,}073)          & 0{,}833 (\(\pm\) 0{,}068)          & 0{,}789 (\(\pm\) 0{,}043)          \\
      Router--specjalista & 0{,}733 (\(\pm\) 0{,}081)          & 0{,}800 (\(\pm\) 0{,}073)          & \textbf{0{,}900 (\(\pm\) 0{,}055)} & \textbf{0{,}811 (\(\pm\) 0{,}041)} \\
      Tablicowa           & \textbf{0{,}833 (\(\pm\) 0{,}068)} & 0{,}767 (\(\pm\) 0{,}077)          & 0{,}767 (\(\pm\) 0{,}077)          & 0{,}789 (\(\pm\) 0{,}043)          \\
      Hierarchiczna       & 0{,}552 (\(\pm\) 0{,}092)          & \textbf{0{,}833 (\(\pm\) 0{,}068)} & 0{,}700 (\(\pm\) 0{,}084)          & 0{,}697 (\(\pm\) 0{,}049)          \\
      ReAct               & 0{,}633 (\(\pm\) 0{,}088)          & 0{,}733 (\(\pm\) 0{,}081)          & 0{,}467 (\(\pm\) 0{,}091)          & 0{,}611 (\(\pm\) 0{,}051)          \\
      \bottomrule
    \end{tabular}
    \zrodlo{opracowanie własne}
  \end{table}

  \subsection{Poprawność}\label{subsec:wyniki-poprawnosc}

  Tabele~\ref{tab:poprawnosc-trudnosc}--\ref{tab:poprawnosc-ubezpieczyciel}
  przedstawiają średnią poprawność odpowiedzi wraz z~odchyleniem standardowym
  w~tym samym rozbiciu. Najwyższą wartość łączną odnotowano dla architektury
  deterministycznej (0{,}716), a~najniższą dla hierarchicznej (0{,}588).
  W~rozbiciu według poziomu trudności poprawność przyjmuje najwyższe wartości dla
  pytań prostych we wszystkich wariantach, natomiast w~przypadku pytań bardzo
  złożonych ulega obniżeniu. W~rozbiciu według kategorii produktowej architektura
  deterministyczna osiągnęła najwyższe wartości we wszystkich trzech grupach.
  W~rozbiciu według ubezpieczyciela nie zaobserwowano jednoznacznej tendencji
  w~wynikach poszczególnych architektur.

  \begin{table}[H]
    \caption{Poprawność według poziomu trudności}\label{tab:poprawnosc-trudnosc}
    \centering
    \footnotesize
    \renewcommand{\arraystretch}{1.3}
    \begin{tabular}{@{} l c c c c c @{}}
      \toprule
      \textbf{Architektura}
      & \textbf{B.~proste}
      & \textbf{Proste}
      & \textbf{Złożone}
      & \textbf{B.~złożone}
      & \textbf{Łącznie}                                                                                                                                                                  \\
      \midrule
      Deterministyczna    & \textbf{0{,}644 (\(\pm\) 0{,}132)} & \textbf{0{,}841 (\(\pm\) 0{,}055)} & \textbf{0{,}685 (\(\pm\) 0{,}077)} & \textbf{0{,}644 (\(\pm\) 0{,}067)} & \textbf{0{,}716 (\(\pm\) 0{,}038)} \\
      Planer--wykonawca   & 0{,}578 (\(\pm\) 0{,}114)          & 0{,}681 (\(\pm\) 0{,}070)          & 0{,}596 (\(\pm\) 0{,}084)          & 0{,}596 (\(\pm\) 0{,}077)          & 0{,}620 (\(\pm\) 0{,}042)          \\
      Router--specjalista & 0{,}611 (\(\pm\) 0{,}116)          & 0{,}830 (\(\pm\) 0{,}053)          & 0{,}633 (\(\pm\) 0{,}078)          & 0{,}544 (\(\pm\) 0{,}073)          & 0{,}663 (\(\pm\) 0{,}039)          \\
      Tablicowa           & 0{,}544 (\(\pm\) 0{,}136)          & 0{,}807 (\(\pm\) 0{,}058)          & 0{,}633 (\(\pm\) 0{,}079)          & 0{,}489 (\(\pm\) 0{,}076)          & 0{,}633 (\(\pm\) 0{,}042)          \\
      Hierarchiczna       & 0{,}567 (\(\pm\) 0{,}119)          & 0{,}678 (\(\pm\) 0{,}076)          & 0{,}608 (\(\pm\) 0{,}077)          & 0{,}485 (\(\pm\) 0{,}079)          & 0{,}588 (\(\pm\) 0{,}043)          \\
      ReAct               & 0{,}633 (\(\pm\) 0{,}105)          & 0{,}704 (\(\pm\) 0{,}077)          & 0{,}567 (\(\pm\) 0{,}081)          & 0{,}515 (\(\pm\) 0{,}082)          & 0{,}599 (\(\pm\) 0{,}044)          \\
      \bottomrule
    \end{tabular}
    \zrodlo{opracowanie własne}
  \end{table}

  \begin{table}[H]
    \caption{Poprawność według kategorii produktowej}\label{tab:poprawnosc-kategoria}
    \centering
    \small
    \renewcommand{\arraystretch}{1.3}
    \begin{tabular}{@{} l c c c c @{}}
      \toprule
      \textbf{Architektura}
      & \textbf{Auto}
      & \textbf{Mieszkanie}
      & \textbf{Podróż}
      & \textbf{Łącznie}                                                                                                                              \\
      \midrule
      Deterministyczna    & \textbf{0{,}700 (\(\pm\) 0{,}068)} & 0{,}683 (\(\pm\) 0{,}076)          & \textbf{0{,}763 (\(\pm\) 0{,}052)} & \textbf{0{,}716 (\(\pm\) 0{,}038)} \\
      Planer--wykonawca   & 0{,}547 (\(\pm\) 0{,}079)          & 0{,}660 (\(\pm\) 0{,}064)          & 0{,}653 (\(\pm\) 0{,}072)          & 0{,}620 (\(\pm\) 0{,}042)          \\
      Router--specjalista & 0{,}593 (\(\pm\) 0{,}071)          & \textbf{0{,}690 (\(\pm\) 0{,}070)} & 0{,}707 (\(\pm\) 0{,}062)          & 0{,}663 (\(\pm\) 0{,}039)          \\
      Tablicowa           & 0{,}667 (\(\pm\) 0{,}066)          & 0{,}603 (\(\pm\) 0{,}077)          & 0{,}630 (\(\pm\) 0{,}074)          & 0{,}633 (\(\pm\) 0{,}042)          \\
      Hierarchiczna       & 0{,}562 (\(\pm\) 0{,}077)          & 0{,}570 (\(\pm\) 0{,}071)          & 0{,}630 (\(\pm\) 0{,}074)          & 0{,}588 (\(\pm\) 0{,}043)          \\
      ReAct               & 0{,}630 (\(\pm\) 0{,}078)          & 0{,}463 (\(\pm\) 0{,}077)          & 0{,}703 (\(\pm\) 0{,}065)          & 0{,}599 (\(\pm\) 0{,}044)          \\
      \bottomrule
    \end{tabular}
    \zrodlo{opracowanie własne}
  \end{table}

  \begin{table}[H]
    \caption{Poprawność według ubezpieczyciela}\label{tab:poprawnosc-ubezpieczyciel}
    \centering
    \small
    \renewcommand{\arraystretch}{1.3}
    \begin{tabular}{@{} l c c c c @{}}
      \toprule
      \textbf{Architektura}
      & \textbf{Hestia}
      & \textbf{PZU}
      & \textbf{Warta}
      & \textbf{Łącznie}                                                                                                                              \\
      \midrule
      Deterministyczna    & 0{,}693 (\(\pm\) 0{,}072)          & \textbf{0{,}793 (\(\pm\) 0{,}057)} & \textbf{0{,}660 (\(\pm\) 0{,}066)} & \textbf{0{,}716 (\(\pm\) 0{,}038)} \\
      Planer--wykonawca   & 0{,}697 (\(\pm\) 0{,}069)          & 0{,}640 (\(\pm\) 0{,}066)          & 0{,}523 (\(\pm\) 0{,}078)          & 0{,}620 (\(\pm\) 0{,}042)          \\
      Router--specjalista & 0{,}673 (\(\pm\) 0{,}070)          & 0{,}750 (\(\pm\) 0{,}060)          & 0{,}567 (\(\pm\) 0{,}070)          & 0{,}663 (\(\pm\) 0{,}039)          \\
      Tablicowa           & 0{,}683 (\(\pm\) 0{,}076)          & 0{,}623 (\(\pm\) 0{,}073)          & 0{,}593 (\(\pm\) 0{,}068)          & 0{,}633 (\(\pm\) 0{,}042)          \\
      Hierarchiczna       & 0{,}603 (\(\pm\) 0{,}078)          & 0{,}707 (\(\pm\) 0{,}065)          & 0{,}453 (\(\pm\) 0{,}072)          & 0{,}588 (\(\pm\) 0{,}043)          \\
      ReAct               & \textbf{0{,}720 (\(\pm\) 0{,}072)} & 0{,}567 (\(\pm\) 0{,}080)          & 0{,}510 (\(\pm\) 0{,}070)          & 0{,}599 (\(\pm\) 0{,}044)          \\
      \bottomrule
    \end{tabular}
    \zrodlo{opracowanie własne}
  \end{table}

  Szczegółowa analiza wyników poprawności w~zestawieniu z~wynikami trafności
  dokumentów (tabela~\ref{tab:trafnosc-ubezpieczyciel}) uwidacznia interesujący
  paradoks. Trafność pobierania dla dokumentów ubezpieczyciela Hestia jest
  systematycznie wyższa (od 0{,}406 do 0{,}480) niż dla dokumentów PZU
  (od 0{,}270 do 0{,}342). Jednakże w~przypadku wiodących architektur to właśnie
  pytania dotyczące PZU osiągają najwyższą poprawność w~fazie syntezy (0{,}793
  w~architekturze deterministycznej), podczas gdy dla Hestii wartości te są
  niższe (0{,}693). Sugeruje to, że dokumenty Hestii charakteryzują się formą
  sprzyjającą wyszukiwaniu wektorowemu (być może lepszą strukturą pojęciową),
  lecz trudniejszą do interpretacji przez model LLM, podczas gdy dokumenty PZU,
  mimo trudniejszego pozycjonowania w~przestrzeni osadzeń, posiadają zapisy,
  z~których łatwiej wydedukować jednoznaczną i~poprawną odpowiedź.

  \subsection{Zwięzłość}\label{subsec:wyniki-zwiezlosc}

  Tabele~\ref{tab:zwiezlosc-trudnosc}--\ref{tab:zwiezlosc-ubezpieczyciel}
  przedstawiają średnią zwięzłość odpowiedzi wraz z~odchyleniem standardowym.
  Wartości są bardziej zbliżone do siebie pomiędzy architekturami niż w~przypadku
  wierności i~poprawności -- mieszczą się w~przedziale 0{,}631--0{,}682.
  Najwyższą wartość łączną odnotowano dla wariantu ReAct (0{,}682), a~najniższą
  dla architektury deterministycznej (0{,}631). W~rozbiciu według poziomu
  trudności wartości zwięzłości nie wykazują jednoznacznej tendencji. W~rozbiciu
  według kategorii i~ubezpieczyciela różnice między architekturami pozostają
  niewielkie.

  \begin{table}[H]
    \caption{Zwięzłość według poziomu trudności}\label{tab:zwiezlosc-trudnosc}
    \centering
    \footnotesize
    \renewcommand{\arraystretch}{1.3}
    \begin{tabular}{@{} l c c c c c @{}}
      \toprule
      \textbf{Architektura}
      & \textbf{B.~proste}
      & \textbf{Proste}
      & \textbf{Złożone}
      & \textbf{B.~złożone}
      & \textbf{Łącznie}                                                                                                                                                                  \\
      \midrule
      Deterministyczna    & 0{,}544 (\(\pm\) 0{,}050)          & 0{,}630 (\(\pm\) 0{,}033)          & 0{,}648 (\(\pm\) 0{,}039)          & 0{,}644 (\(\pm\) 0{,}033)          & 0{,}631 (\(\pm\) 0{,}019)          \\
      Planer--wykonawca   & 0{,}667 (\(\pm\) 0{,}050)          & 0{,}637 (\(\pm\) 0{,}039)          & 0{,}630 (\(\pm\) 0{,}037)          & 0{,}689 (\(\pm\) 0{,}038)          & 0{,}653 (\(\pm\) 0{,}021)          \\
      Router--specjalista & \textbf{0{,}700 (\(\pm\) 0{,}047)} & 0{,}659 (\(\pm\) 0{,}031)          & 0{,}674 (\(\pm\) 0{,}035)          & 0{,}581 (\(\pm\) 0{,}037)          & 0{,}644 (\(\pm\) 0{,}019)          \\
      Tablicowa           & 0{,}678 (\(\pm\) 0{,}060)          & 0{,}674 (\(\pm\) 0{,}035)          & 0{,}626 (\(\pm\) 0{,}039)          & 0{,}652 (\(\pm\) 0{,}039)          & 0{,}653 (\(\pm\) 0{,}021)          \\
      Hierarchiczna       & 0{,}700 (\(\pm\) 0{,}047)          & \textbf{0{,}700 (\(\pm\) 0{,}038)} & \textbf{0{,}700 (\(\pm\) 0{,}037)} & 0{,}622 (\(\pm\) 0{,}040)          & 0{,}676 (\(\pm\) 0{,}021)          \\
      ReAct               & 0{,}611 (\(\pm\) 0{,}060)          & 0{,}685 (\(\pm\) 0{,}035)          & 0{,}689 (\(\pm\) 0{,}034)          & \textbf{0{,}696 (\(\pm\) 0{,}041)} & \textbf{0{,}682 (\(\pm\) 0{,}020)} \\
      \bottomrule
    \end{tabular}
    \zrodlo{opracowanie własne}
  \end{table}

  \begin{table}[H]
    \caption{Zwięzłość według kategorii produktowej}\label{tab:zwiezlosc-kategoria}
    \centering
    \small
    \renewcommand{\arraystretch}{1.3}
    \begin{tabular}{@{} l c c c c @{}}
      \toprule
      \textbf{Architektura}
      & \textbf{Auto}
      & \textbf{Mieszkanie}
      & \textbf{Podróż}
      & \textbf{Łącznie}                                                                                                                              \\
      \midrule
      Deterministyczna    & 0{,}610 (\(\pm\) 0{,}035)          & 0{,}630 (\(\pm\) 0{,}034)          & 0{,}653 (\(\pm\) 0{,}029)          & 0{,}631 (\(\pm\) 0{,}019)          \\
      Planer--wykonawca   & 0{,}700 (\(\pm\) 0{,}038)          & 0{,}613 (\(\pm\) 0{,}032)          & 0{,}647 (\(\pm\) 0{,}035)          & 0{,}653 (\(\pm\) 0{,}021)          \\
      Router--specjalista & 0{,}617 (\(\pm\) 0{,}034)          & \textbf{0{,}673 (\(\pm\) 0{,}031)} & 0{,}643 (\(\pm\) 0{,}033)          & 0{,}644 (\(\pm\) 0{,}019)          \\
      Tablicowa           & 0{,}657 (\(\pm\) 0{,}035)          & 0{,}623 (\(\pm\) 0{,}037)          & 0{,}680 (\(\pm\) 0{,}034)          & 0{,}653 (\(\pm\) 0{,}021)          \\
      Hierarchiczna       & 0{,}679 (\(\pm\) 0{,}031)          & 0{,}640 (\(\pm\) 0{,}038)          & \textbf{0{,}710 (\(\pm\) 0{,}038)} & 0{,}676 (\(\pm\) 0{,}021)          \\
      ReAct               & \textbf{0{,}737 (\(\pm\) 0{,}031)} & 0{,}660 (\(\pm\) 0{,}039)          & 0{,}650 (\(\pm\) 0{,}033)          & \textbf{0{,}682 (\(\pm\) 0{,}020)} \\
      \bottomrule
    \end{tabular}
    \zrodlo{opracowanie własne}
  \end{table}

  \begin{table}[H]
    \caption{Zwięzłość według ubezpieczyciela}\label{tab:zwiezlosc-ubezpieczyciel}
    \centering
    \small
    \renewcommand{\arraystretch}{1.3}
    \begin{tabular}{@{} l c c c c @{}}
      \toprule
      \textbf{Architektura}
      & \textbf{Hestia}
      & \textbf{PZU}
      & \textbf{Warta}
      & \textbf{Łącznie}                                                                                                                              \\
      \midrule
      Deterministyczna    & 0{,}647 (\(\pm\) 0{,}031)          & 0{,}597 (\(\pm\) 0{,}033)          & 0{,}650 (\(\pm\) 0{,}035)          & 0{,}631 (\(\pm\) 0{,}019)          \\
      Planer--wykonawca   & 0{,}633 (\(\pm\) 0{,}033)          & 0{,}673 (\(\pm\) 0{,}039)          & 0{,}653 (\(\pm\) 0{,}035)          & 0{,}653 (\(\pm\) 0{,}021)          \\
      Router--specjalista & 0{,}660 (\(\pm\) 0{,}027)          & 0{,}640 (\(\pm\) 0{,}031)          & 0{,}633 (\(\pm\) 0{,}039)          & 0{,}644 (\(\pm\) 0{,}019)          \\
      Tablicowa           & \textbf{0{,}697 (\(\pm\) 0{,}035)} & 0{,}653 (\(\pm\) 0{,}036)          & 0{,}610 (\(\pm\) 0{,}034)          & 0{,}653 (\(\pm\) 0{,}021)          \\
      Hierarchiczna       & 0{,}690 (\(\pm\) 0{,}035)          & 0{,}670 (\(\pm\) 0{,}033)          & 0{,}670 (\(\pm\) 0{,}040)          & 0{,}676 (\(\pm\) 0{,}021)          \\
      ReAct               & 0{,}647 (\(\pm\) 0{,}031)          & \textbf{0{,}697 (\(\pm\) 0{,}040)} & \textbf{0{,}703 (\(\pm\) 0{,}033)} & \textbf{0{,}682 (\(\pm\) 0{,}020)} \\
      \bottomrule
    \end{tabular}
    \zrodlo{opracowanie własne}
  \end{table}

  \subsection{Trafność dokumentów}\label{subsec:wyniki-trafnosc}

  Trafność dokumentów oceniano na poziomie pojedynczego cytowania -- łącznie
  1\,786~par pytanie--fragment.
  Tabele~\ref{tab:trafnosc-trudnosc}--\ref{tab:trafnosc-ubezpieczyciel} podają
  średnie arytmetyczne ocen binarnych wraz z~odchyleniem standardowym w~obrębie
  odpowiednich podzbiorów. Zaobserwowane wartości trafności są niższe niż
  w~przypadku pozostałych miar jakościowych i~mieszczą się w~przedziale
  0{,}332--0{,}364 dla wyników łącznych. Architektura deterministyczna i~ReAct
  osiągnęły tę samą najwyższą wartość (0{,}364), a~tablicowa najniższą (0{,}332).
  We wszystkich architekturach trafność przyjmuje wyższe wartości dla pytań
  bardzo prostych i~maleje wraz ze wzrostem poziomu trudności. W~rozbiciu według
  ubezpieczyciela najwyższe wartości we wszystkich architekturach odnotowano
  w~grupie \texttt{Hestia}, a~najniższe -- w~grupie \texttt{Warta}.

  \begin{table}[H]
    \caption{Trafność dokumentów według poziomu trudności}\label{tab:trafnosc-trudnosc}
    \centering
    \footnotesize
    \renewcommand{\arraystretch}{1.3}
    \begin{tabular}{@{} l c c c c c @{}}
      \toprule
      \textbf{Architektura}
      & \textbf{B.~proste}
      & \textbf{Proste}
      & \textbf{Złożone}
      & \textbf{B.~złożone}
      & \textbf{Łącznie}                                                                                                                                                                  \\
      \midrule
      Deterministyczna    & \textbf{0{,}545 (\(\pm\) 0{,}087)} & \textbf{0{,}394 (\(\pm\) 0{,}048)} & 0{,}330 (\(\pm\) 0{,}046)          & 0{,}305 (\(\pm\) 0{,}047)          & \textbf{0{,}364 (\(\pm\) 0{,}026)} \\
      Planer--wykonawca   & 0{,}500 (\(\pm\) 0{,}091)          & 0{,}323 (\(\pm\) 0{,}048)          & 0{,}359 (\(\pm\) 0{,}050)          & 0{,}315 (\(\pm\) 0{,}049)          & 0{,}349 (\(\pm\) 0{,}027)          \\
      Router--specjalista & 0{,}441 (\(\pm\) 0{,}085)          & 0{,}392 (\(\pm\) 0{,}050)          & 0{,}283 (\(\pm\) 0{,}045)          & \textbf{0{,}370 (\(\pm\) 0{,}048)} & 0{,}358 (\(\pm\) 0{,}026)          \\
      Tablicowa           & 0{,}500 (\(\pm\) 0{,}086)          & 0{,}307 (\(\pm\) 0{,}046)          & 0{,}337 (\(\pm\) 0{,}048)          & 0{,}298 (\(\pm\) 0{,}045)          & 0{,}332 (\(\pm\) 0{,}026)          \\
      Hierarchiczna       & 0{,}515 (\(\pm\) 0{,}087)          & 0{,}330 (\(\pm\) 0{,}046)          & 0{,}333 (\(\pm\) 0{,}047)          & 0{,}314 (\(\pm\) 0{,}046)          & 0{,}344 (\(\pm\) 0{,}026)          \\
      ReAct               & 0{,}500 (\(\pm\) 0{,}134)          & 0{,}370 (\(\pm\) 0{,}066)          & \textbf{0{,}419 (\(\pm\) 0{,}075)} & 0{,}219 (\(\pm\) 0{,}073)          & 0{,}364 (\(\pm\) 0{,}040)          \\
      \bottomrule
    \end{tabular}
    \zrodlo{opracowanie własne}
  \end{table}

  \begin{table}[H]
    \caption{Trafność dokumentów według kategorii produktowej}\label{tab:trafnosc-kategoria}
    \centering
    \small
    \renewcommand{\arraystretch}{1.3}
    \begin{tabular}{@{} l c c c c @{}}
      \toprule
      \textbf{Architektura}
      & \textbf{Auto}
      & \textbf{Mieszkanie}
      & \textbf{Podróż}
      & \textbf{Łącznie}                                                                                                                              \\
      \midrule
      Deterministyczna    & 0{,}315 (\(\pm\) 0{,}044)          & \textbf{0{,}386 (\(\pm\) 0{,}046)} & 0{,}391 (\(\pm\) 0{,}047)          & \textbf{0{,}364 (\(\pm\) 0{,}026)} \\
      Planer--wykonawca   & 0{,}323 (\(\pm\) 0{,}048)          & 0{,}383 (\(\pm\) 0{,}047)          & 0{,}337 (\(\pm\) 0{,}047)          & 0{,}349 (\(\pm\) 0{,}027)          \\
      Router--specjalista & \textbf{0{,}352 (\(\pm\) 0{,}046)} & 0{,}372 (\(\pm\) 0{,}045)          & 0{,}349 (\(\pm\) 0{,}046)          & 0{,}358 (\(\pm\) 0{,}026)          \\
      Tablicowa           & 0{,}315 (\(\pm\) 0{,}044)          & 0{,}342 (\(\pm\) 0{,}044)          & 0{,}339 (\(\pm\) 0{,}045)          & 0{,}332 (\(\pm\) 0{,}026)          \\
      Hierarchiczna       & 0{,}309 (\(\pm\) 0{,}044)          & 0{,}331 (\(\pm\) 0{,}043)          & \textbf{0{,}394 (\(\pm\) 0{,}047)} & 0{,}344 (\(\pm\) 0{,}026)          \\
      ReAct               & 0{,}318 (\(\pm\) 0{,}070)          & 0{,}372 (\(\pm\) 0{,}074)          & 0{,}393 (\(\pm\) 0{,}065)          & 0{,}364 (\(\pm\) 0{,}040)          \\
      \bottomrule
    \end{tabular}
    \zrodlo{opracowanie własne}
  \end{table}

  \begin{table}[H]
    \caption{Trafność dokumentów według ubezpieczyciela}\label{tab:trafnosc-ubezpieczyciel}
    \centering
    \small
    \renewcommand{\arraystretch}{1.3}
    \begin{tabular}{@{} l c c c c @{}}
      \toprule
      \textbf{Architektura}
      & \textbf{Hestia}
      & \textbf{PZU}
      & \textbf{Warta}
      & \textbf{Łącznie}                                                                                                                              \\
      \midrule
      Deterministyczna    & 0{,}460 (\(\pm\) 0{,}047)          & 0{,}336 (\(\pm\) 0{,}044)          & 0{,}294 (\(\pm\) 0{,}044)          & \textbf{0{,}364 (\(\pm\) 0{,}026)} \\
      Planer--wykonawca   & 0{,}406 (\(\pm\) 0{,}048)          & 0{,}333 (\(\pm\) 0{,}048)          & 0{,}304 (\(\pm\) 0{,}046)          & 0{,}349 (\(\pm\) 0{,}027)          \\
      Router--specjalista & 0{,}436 (\(\pm\) 0{,}047)          & 0{,}324 (\(\pm\) 0{,}046)          & \textbf{0{,}313 (\(\pm\) 0{,}043)} & 0{,}358 (\(\pm\) 0{,}026)          \\
      Tablicowa           & 0{,}425 (\(\pm\) 0{,}047)          & 0{,}270 (\(\pm\) 0{,}042)          & 0{,}301 (\(\pm\) 0{,}043)          & 0{,}332 (\(\pm\) 0{,}026)          \\
      Hierarchiczna       & 0{,}416 (\(\pm\) 0{,}046)          & \textbf{0{,}342 (\(\pm\) 0{,}045)} & 0{,}274 (\(\pm\) 0{,}042)          & 0{,}344 (\(\pm\) 0{,}026)          \\
      ReAct               & \textbf{0{,}480 (\(\pm\) 0{,}071)} & 0{,}302 (\(\pm\) 0{,}063)          & 0{,}300 (\(\pm\) 0{,}072)          & 0{,}364 (\(\pm\) 0{,}040)          \\
      \bottomrule
    \end{tabular}
    \zrodlo{opracowanie własne}
  \end{table}

  \section{Parametry konfiguracyjne eksperymentu}\label{sec:parametry-konfiguracyjne}

  Sekcja dokumentuje wartości parametrów sterujących przebiegiem uruchomienia.
  Pełną konfigurację zawierają moduły \texttt{sources/config/app.py}
  i~\texttt{sources/config/graph.py}; poniższe zestawienie obejmuje wyłącznie
  ustawienia użyte w~uruchomieniu \texttt{20260313\_200511}.

  \subsection{Parametry modeli językowych}\label{subsec:parametry-modeli}

  Wszystkie role pomocnicze -- parser, router, planer, kontroler ReAct, rewriter,
  dekompozytor i~ewaluator -- zrealizowano z~wykorzystaniem modelu
  \texttt{gpt-5}. Ten sam model zastosowano do syntezy odpowiedzi końcowej. Każdą
  rolę skonfigurowano z~temperaturą równą zero, limitem 120~s na pojedyncze
  wywołanie i~mechanizmem ponowień z~wykładniczym opóźnieniem od 1 do 60~s.
  Tabela~\ref{tab:parametry-modeli} zestawia pełną konfigurację.

  \begin{table}[H]
    \caption{Parametry modeli językowych}\label{tab:parametry-modeli}
    \centering
    \small
    \renewcommand{\arraystretch}{1.3}
    \begin{tabular}{@{} l l l @{}}
      \toprule
      \textbf{Parametr}
      & \textbf{Wartość}
      & \textbf{Rola}                                               \\
      \midrule
      \texttt{model\_name}
      & \texttt{gpt-5}   & model syntezy odpowiedzi                 \\
      \texttt{parser\_model\_name}
      & \texttt{gpt-5}   & parser pytań                             \\
      \texttt{router\_model\_name}
      & \texttt{gpt-5}   & router specjalisty                       \\
      \texttt{planner\_model\_name}
      & \texttt{gpt-5}   & planer kroków                            \\
      \texttt{controller\_model\_name}
      & \texttt{gpt-5}   & kontroler ReAct                          \\
      \texttt{rewrite\_model\_name}
      & \texttt{gpt-5}   & przepisywanie zapytań                    \\
      \texttt{decomposer\_model\_name}
      & \texttt{gpt-5}   & dekompozycja pytań                       \\
      \texttt{evaluation\_model\_name}
      & \texttt{gpt-5}   & sędzia jakości                           \\
      \texttt{temperature}
      & 0{,}0             & determinizm wyjścia                      \\
      \texttt{retry\_max\_attempts}
      & 5                & liczba ponowień                          \\
      \texttt{retry\_min\_wait}
      & 1{,}0~s          & minimalna pauza między ponowieniami      \\
      \texttt{retry\_max\_wait}
      & 60{,}0~s         & maksymalna pauza między ponowieniami     \\
      \texttt{llm\_request\_timeout}
      & 120{,}0~s        & limit pojedynczego wywołania             \\
      \texttt{max\_concurrent\_api\_calls}
      & 16               & limit równoległości wywołań \texttt{API} \\
      \bottomrule
    \end{tabular}
    \zrodlo{opracowanie własne na podstawie kodu źródłowego}
  \end{table}

  \subsection{Parametry wyszukiwania}\label{subsec:parametry-wyszukiwania}

  Wyszukiwanie wektorowe wykonano na kolekcji \texttt{insurance\_documents}
  w~bazie \texttt{Qdrant}, z~wykorzystaniem modelu osadzeń
  \texttt{text-embedding-3-large} o~wymiarze 3072. Etap rerankingu wyłączono we
  wszystkich architekturach; węzeł rerankera pozostał jednak obecny w~każdym
  grafie. Tabela~\ref{tab:parametry-wyszukiwania} prezentuje pełną konfigurację
  warstwy wyszukiwania wraz z~nominalnymi parametrami rerankera.

  \begin{table}[H]
    \caption{Parametry wyszukiwania}\label{tab:parametry-wyszukiwania}
    \centering
    \small
    \renewcommand{\arraystretch}{1.3}
    \begin{tabularx}{\textwidth}{@{} l l X @{}}
      \toprule
      \textbf{Parametr}
      & \textbf{Wartość}
      & \textbf{Opis}                                                              \\
      \midrule
      \texttt{embedding.model\_name}
      & \texttt{text-embedding-3-large} & model osadzeń                            \\
      \texttt{embedding.dimension}
      & 3072                            & wymiar wektora                           \\
      \texttt{qdrant.collection\_name}
      & \texttt{insurance\_documents}   & nazwa kolekcji                           \\
      \texttt{top\_k}
      & 5                               & liczba pobieranych fragmentów            \\
      \texttt{evidence\_top\_k}
      & 4                               & liczba dowodów przekazywanych do syntezy \\
      \texttt{reranking.enabled}
      & \texttt{False}                  & reranking wyłączony w~eksperymencie      \\
      \texttt{reranking.score\_threshold}
      & 0{,}3                            & próg odrzucenia fragmentu                \\
      \texttt{reranking.top\_k\_before}
      & 20                              & wejście do rerankera                     \\
      \texttt{reranking.top\_k\_after}
      & 5                               & wyjście z~rerankera                      \\
      \bottomrule
    \end{tabularx}
    \zrodlo{opracowanie własne na podstawie kodu źródłowego}
  \end{table}

  \subsection{Parametry fragmentacji}\label{subsec:parametry-fragmentacji}

  Dokumenty OWU i~IPID przetwarzano przy użyciu odrębnych profili fragmentacji.
  W~obu przypadkach zastosowano podział po nagłówkach Markdown poziomów 1--3,
  a~następnie rekurencyjną fragmentację większych sekcji. Parametry wyrażone
  w~liczbie znaków zestawiono w~tabeli~\ref{tab:parametry-fragmentacji};
  zastosowanie krótszych fragmentów dla IPID wynikało z~większej gęstości
  informacyjnej tego formatu.

  \begin{table}[H]
    \caption{Parametry fragmentacji dokumentów}\label{tab:parametry-fragmentacji}
    \centering
    \small
    \renewcommand{\arraystretch}{1.3}
    \begin{tabular}{@{} l c c l @{}}
      \toprule
      \textbf{Parametr}
      & \textbf{OWU}
      & \textbf{IPID}
      & \textbf{Opis}                                          \\
      \midrule
      \texttt{chunk\_size}      & 512           & 384 & docelowy rozmiar fragmentu       \\
      \texttt{chunk\_overlap}   & 64            & 48  & nakładanie sąsiednich fragmentów \\
      \texttt{min\_chunk\_size} & 64            & 48  & minimalny dopuszczalny rozmiar   \\
      \bottomrule
    \end{tabular}
    \zrodlo{opracowanie własne na podstawie kodu źródłowego}
  \end{table}

  \subsection{Parametry architektur}\label{subsec:parametry-architektur}

  Architektury iteracyjne wykorzystują jawne ograniczenia liczby kroków;
  architektura hierarchiczna ogranicza dekompozycję do trzech podpytań
  bezpośrednio w~kodzie węzła dekomponera. Architektury deterministyczna
  i~router--specjalista nie mają parametrów sterujących właściwych dla danej
  topologii. Tabela~\ref{tab:parametry-architektur} zestawia wartości dla
  wariantów, w~których takie parametry występują.

  \begin{table}[H]
    \caption{Parametry sterujące architektur}\label{tab:parametry-architektur}
    \centering
    \small
    \renewcommand{\arraystretch}{1.3}
    \begin{tabular}{@{} l l c l @{}}
      \toprule
      \textbf{Architektura}
      & \textbf{Parametr}
      & \textbf{Wartość}
      & \textbf{Znaczenie}                            \\
      \midrule
      ReAct
      & \texttt{max\_react\_iterations}          & 15
      & limit iteracji pętli obserwacyjnej            \\
      Tablicowa
      & \texttt{max\_blackboard\_iterations}     & 15
      & limit cykli dyspozytora                       \\
      Hierarchiczna
      & max. liczba podpytań                     & 3
      & przycięcie listy \texttt{sub\_questions}      \\
      \bottomrule
    \end{tabular}
    \zrodlo{opracowanie własne na podstawie kodu źródłowego}
  \end{table}

  \subsection{Parametry wykonania}\label{subsec:parametry-wykonania}

  Komponent wykonawczy pobierał parametry współbieżności
  z~\texttt{ConcurrencyConfig} i~inicjował pulę 32~wątków obsługujących pełną
  macierz par pytanie--architektura. Tabela~\ref{tab:parametry-wykonania}
  przedstawia liczebności i~ustawienia dla uruchomienia \texttt{20260313\_200511}.

  \begin{table}[H]
    \caption{Parametry wykonania eksperymentu}\label{tab:parametry-wykonania}
    \centering
    \small
    \renewcommand{\arraystretch}{1.3}
    \begin{tabular}{@{} l c l @{}}
      \toprule
      \textbf{Parametr}
      & \textbf{Wartość}
      & \textbf{Opis}                                            \\
      \midrule
      \texttt{max\_workers}    & 32               & rozmiar puli wątków roboczych         \\
      Liczba pytań testowych   & 90               & zestaw \texttt{pytania\_20251130.csv} \\
      Liczba architektur       & 6                & komplet \texttt{ALL\_PATTERNS}        \\
      Łączna liczba uruchomień & 540              & macierz \(90 \times 6\)               \\
      \bottomrule
    \end{tabular}
    \zrodlo{opracowanie własne na podstawie kodu źródłowego}
  \end{table}